\section{Advanced Detection Patterns}

Advanced detection engineering focuses on identifying subtle, high‑value behaviors that adversaries rely on to evade defenses, escalate privileges, and maintain persistence. These patterns go beyond simple indicators or basic anomalies; they capture the underlying techniques that remain consistent even as tooling, infrastructure, and malware families evolve. This section introduces several advanced detection patterns, each supported by SPL and KQL logic, along with Linux investigation pivots. The goal is to provide a deep understanding of how these behaviors manifest in telemetry and how to detect them reliably in enterprise environments.

\subsection*{Detection Pattern: Parent PID Spoofing}

\textbf{Behavior and Rationale}

Parent PID spoofing occurs when an adversary manipulates the parent process ID of a newly created process to make it appear as though it was launched by a legitimate parent. This technique is used to evade behavioral detections that rely on parent-child relationships. While Windows does not natively allow arbitrary parent PID assignment, certain APIs and tools (such as process hollowing or direct system calls) can create misleading process trees. Detecting this behavior requires identifying inconsistencies between expected and observed parent-child relationships.

\textbf{SPL Detection Logic}
\begin{lstlisting}
index=sysmon EventCode=1
| eval expected_parent=case(
    Image="*powershell.exe", "explorer.exe",
    Image="*cmd.exe", "explorer.exe",
    Image="*wscript.exe", "explorer.exe",
    1==1, "unknown")
| where NOT like(ParentImage, "%" . expected_parent . "%")
| table _time, Computer, User, ParentImage, Image, CommandLine
\end{lstlisting}

\textbf{KQL Detection Logic}
\begin{lstlisting}
Sysmon
| where EventID == 1
| extend ExpectedParent = case(
    Image endswith "powershell.exe", "explorer.exe",
    Image endswith "cmd.exe", "explorer.exe",
    Image endswith "wscript.exe", "explorer.exe",
    "unknown")
| where ParentImage !contains ExpectedParent
| project TimeGenerated, Computer, User, ParentImage, Image, CommandLine
\end{lstlisting}

\textbf{Linux Pivot}
\begin{lstlisting}[language=bash]
ps -eo pid,ppid,cmd --forest
\end{lstlisting}

\textbf{Explanation}

This detection identifies processes whose parent does not match expected norms. While not all mismatches are malicious, unexpected parent-child relationships are a strong indicator of process injection, hollowing, or spoofing.

\subsection*{Detection Pattern: LOLBins Abuse}

\textbf{Behavior and Rationale}

Living‑off‑the‑land binaries (LOLBins) are legitimate system tools that adversaries abuse to execute code, download payloads, or evade detection. Examples include \texttt{certutil}, \texttt{mshta}, \texttt{wmic}, and \texttt{rundll32}. Because these tools are signed and trusted, traditional signature-based detections often fail. Behavioral detections must focus on unusual usage patterns, command-line arguments, and execution context.

\textbf{SPL Detection Logic}
\begin{lstlisting}
index=sysmon EventCode=1
| search Image IN ("*certutil.exe", "*mshta.exe", "*wmic.exe", "*rundll32.exe")
| table _time, Computer, User, ParentImage, Image, CommandLine
\end{lstlisting}

\textbf{KQL Detection Logic}
\begin{lstlisting}
Sysmon
| where EventID == 1
| where Image endswith "certutil.exe"
    or Image endswith "mshta.exe"
    or Image endswith "wmic.exe"
    or Image endswith "rundll32.exe"
| project TimeGenerated, Computer, User, ParentImage, Image, CommandLine
\end{lstlisting}

\textbf{Linux Pivot}
\begin{lstlisting}[language=bash]
grep -E "curl|wget|python|perl" /var/log/syslog
\end{lstlisting}

\textbf{Explanation}

This detection highlights the execution of high‑risk binaries. Analysts should pay close attention to suspicious command-line arguments, such as Base64 strings, remote URLs, or script execution flags.

\subsection*{Detection Pattern: Credential Dumping Behavior}

\textbf{Behavior and Rationale}

Credential dumping is a critical step in many intrusions, enabling adversaries to escalate privileges and move laterally. Tools such as Mimikatz, LSASS memory readers, and custom malware often interact with sensitive processes or registry hives. Detecting credential dumping requires monitoring for access to LSASS, suspicious handle requests, or unusual process memory reads.

\textbf{SPL Detection Logic}
\begin{lstlisting}
index=sysmon EventCode=10 TargetImage="*lsass.exe"
| table _time, Computer, User, SourceImage, TargetImage, GrantedAccess
\end{lstlisting}

\textbf{KQL Detection Logic}
\begin{lstlisting}
Sysmon
| where EventID == 10 and TargetImage endswith "lsass.exe"
| project TimeGenerated, Computer, User, SourceImage, TargetImage, GrantedAccess
\end{lstlisting}

\textbf{Linux Pivot}
\begin{lstlisting}[language=bash]
grep -i "shadow" /var/log/auth.log
\end{lstlisting}

\textbf{Explanation}

Sysmon Event ID 10 captures process access attempts. Any process attempting to access LSASS with high privileges is suspicious and should be investigated immediately.

\subsection*{Detection Pattern: Registry-Based Persistence}

\textbf{Behavior and Rationale}

Adversaries frequently use registry keys to establish persistence, especially via \texttt{Run}, \texttt{RunOnce}, and \texttt{Services}. Detecting modifications to these keys is essential for identifying long-term footholds.

\textbf{SPL Detection Logic}
\begin{lstlisting}
index=sysmon EventCode=13
| search TargetObject="*\\Run*" OR TargetObject="*\\RunOnce*"
| table _time, Computer, User, TargetObject, Details
\end{lstlisting}

\textbf{KQL Detection Logic}
\begin{lstlisting}
Sysmon
| where EventID == 13
| where TargetObject contains "\\Run" or TargetObject contains "\\RunOnce"
| project TimeGenerated, Computer, User, TargetObject, Details
\end{lstlisting}

\textbf{Linux Pivot}
\begin{lstlisting}[language=bash]
ls -al ~/.config/autostart/
\end{lstlisting}

\textbf{Explanation}

Registry persistence is common in malware and should be monitored closely. Analysts should validate whether the modified entry corresponds to legitimate software.

\subsection*{Detection Pattern: Suspicious Service Creation}

\textbf{Behavior and Rationale}

Creating or modifying services is a powerful persistence and privilege escalation technique. Adversaries often create services that execute malicious binaries with SYSTEM privileges.

\textbf{SPL Detection Logic}
\begin{lstlisting}
index=wineventlog EventCode=7045
| table _time, ServiceName, ImagePath, StartType, AccountName
\end{lstlisting}

\textbf{KQL Detection Logic}
\begin{lstlisting}
SecurityEvent
| where EventID == 7045
| project TimeGenerated, ServiceName, ImagePath, StartType, Account
\end{lstlisting}

\textbf{Linux Pivot}
\begin{lstlisting}[language=bash]
systemctl list-units --type=service
\end{lstlisting}

\textbf{Explanation}

Service creation is rare in most environments and should be treated as high-risk unless associated with known administrative activity.

\subsection*{Summary}

The advanced detection patterns in this section focus on behaviors that adversaries rely on to evade defenses and maintain access. By understanding these patterns and implementing the associated SPL and KQL logic, analysts can build high-fidelity detections that remain effective even as threats evolve. These patterns form the foundation of a mature detection engineering program and support long-term operational resilience.
